\documentclass[12pt,a4paper]{article}
\usepackage[utf8]{inputenc}
\usepackage[T2A]{fontenc}
\usepackage[russian]{babel}
\usepackage{amsmath,amssymb,amsthm}
\usepackage{booktabs,array}
\usepackage{geometry}
\usepackage{microtype}
\geometry{margin=2.3cm}

\newtheorem{proposition}{Предложение}
\newtheorem{nogocriterion}{Критерий провала}

\title{Исчерпание минимальных семейных маршрутов:\\
affine incidence против projective carrier}
\author{}
\date{5 августа 2026 г.}

\begin{document}
\maketitle

\section{Область исчерпания}

Настоящий gate не перебирает все мыслимые теории flavour. Он исчерпывает три
минимальных класса, непосредственно продолжающих уже объявленное меню:
\begin{enumerate}
 \item операторы incidence, порождённые всеми элементами $AGL(2,2)$ на
 четырёх spin-состояниях с последующим ограничением на триплет;
 \item действия нетривиального projective $\mathbb Z_2$-flux на носителях
 размерности $2m$;
 \item intrinsic cohomology, product-metric, cellular-Dirac и determinant-line
 selectors на пространстве spin-структур.
\end{enumerate}

\section{Полный affine-incidence перебор}

Текущие две трансляции и shear порождают подгруппу порядка $8$. Её операторная
алгебра на семейном триплете имеет размерность $5$ и коммутант размерности $2$,
то есть сохраняет разложение $1+2$.

В полном $AGL(2,2)\simeq S_4$ остаётся шестнадцать элементов вне этой
подгруппы. Они распадаются на три орбиты текущей геометрической группы:
\begin{center}
\begin{tabular}{ccc}
\toprule
Тип цикла & Число кандидатов & Результат на триплете \\
\midrule
$2+1+1$ & $4$ & полная $M_3(\mathbb C)$ \\
$3+1$ & $8$ & полная $M_3(\mathbb C)$ \\
$4$ & $4$ & редуцируемая алгебра размерности $5$ \\
\bottomrule
\end{tabular}
\end{center}
Таким образом, двенадцать из шестнадцати одиночных дополнительных incidence-
операторов снимают алгебраическое препятствие и допускают полное трёхсемейное
смешивание.

Но это не вывод selector-а. Факторы $\mathbb{RP}^3$ и $S^1$ неэквивалентны,
поэтому недостающие affine-преобразования являются симметриями абстрактного
меню, а не геометрическими автоморфизмами носителя. Усреднение любого кандидата
по текущей группе порядка $8$ возвращает его в коммутант и восстанавливает
$1+2$.

\begin{proposition}[Existence без selection]
Полная $M_3$-алгебра семейного смешивания существует уже в минимальном
четырёхсостоянийном меню. Нерешённой является не задача существования, а выбор
одного из двенадцати рабочих направлений, его коэффициента относительно
факторного лапласиана и его назначения up/down-сектору.
\end{proposition}

\section{Исчерпание projective-carrier}

Нетривиальный flux задаёт
\begin{equation}
 U^2=V^2=I,
 \qquad UV=-VU.
\end{equation}
В четырёх измерениях порождённая алгебра имеет вид
\begin{equation}
 M_2(\mathbb C)\otimes I_2,
\end{equation}
а её коммутант равен $I_2\otimes M_2(\mathbb C)$. Общий эрмитов flux-only
оператор
\begin{equation}
 H=aI+bU+cV+d\,iUV
\end{equation}
имеет спектр
\begin{equation}
 a-r,\ a-r,\ a+r,\ a+r,
 \qquad r=\sqrt{b^2+c^2+d^2}.
\end{equation}
Поэтому его спектральные проекторы имеют только ранги $0,2,4$: split $1+3$
невозможен.

Общий носитель имеет форму
\begin{equation}
 \mathbb C^2_{flux}\otimes\mathbb C^m_{mult}.
\end{equation}
Попытка взять $m=3$ и назвать multiplicity-пространство тремя поколениями
даёт шестимерный носитель, но flux действует на нём тождественно. Вся массовая
иерархия лежит в свободном коммутанте $M_3(\mathbb C)$; две CKM-матрицы требуют
двух некоммутирующих элементов этого коммутанта. Топология их не выбирает.

\begin{nogocriterion}[Провал минимального projective rescue]
Четырёхмерный flux не даёт $1+3$ из-за чётной спектральной кратности.
Шестимерный вариант возвращает число три только как необъяснённую кратность и
переносит всю flavour-задачу в произвольный $M_3$-коммутант. Следовательно,
минимальная projective-ветвь не замыкает generation count и mixing одновременно.
\end{nogocriterion}

\section{Cohomology и metric selector}

Остаётся проверить, не выбирают ли недостающий incidence-оператор сама
cohomology ring или неравные длины факторов. Над $\mathbb F_2$
\begin{equation}
 H^*(\mathbb{RP}^3\times S^1)
 =\mathbb F_2[a,b]/(a^4,b^2).
\end{equation}
Её линейная группа автоморфизмов на $H^1$ содержит только два элемента:
тождество и $a\mapsto a+b$, $b\mapsto b$. Поэтому три ненулевых класса имеют
орбиты
\begin{equation}
 \{b\},\qquad\{a,a+b\},
\end{equation}
а коммутант естественного действия снова имеет размерность $5$ и структуру
$1+2$.

При этом topology и product metric действительно различают все три класса.
Индикатор ненулевого квадрата, наличие свободной $S^1$-компоненты и длины
\begin{equation}
 \ell(a)=\pi,\qquad \ell(b)=2\pi,\qquad
 \ell(a+b)=\sqrt5\,\pi
\end{equation}
дают три различные сигнатуры. Но соответствующие операторы диагональны и
попарно коммутируют. Они могут маркировать поколения и задавать иерархию, но
предсказывают тождественное mixing.

Систолический дефект на $\mathbb{RP}^1$ выбирает класс $a$. Однако его
локализация в фиксированной точке $S^1$ вводит неканоническую фазу/точку.
Оборачивание дефекта вдоль всего $S^1$ удаляет эту свободу, но сохраняет
коммутирующую факторную структуру. Поэтому defect-route не выводит selector без
того же дополнительного boundary datum, которое ранее было признано open.

\section{Единственный живой минимальный маршрут}

После трёх переборов остаётся один конкретный тип положительного исхода:
необходимо вывести из новой предварительной структуры один из двенадцати
outside-$D_8$ incidence-операторов. Допустимыми источниками могут быть только
заранее объявленные boundary data, дискретный Dirac-комплекс или новая
геометрическая связь между двумя факторами. Эта структура должна одновременно
фиксировать:
\begin{enumerate}
 \item конкретную affine-орбиту и элемент внутри неё;
 \item относительный вес к факторному оператору;
 \item различное, но связанное назначение up- и down-Yukawa;
 \item минимум два независимых mixing-направления.
\end{enumerate}

Если такой selector не выводится, семейная ветвь текущего конечного меню
закрывается как физическая модель. При этом математический результат остаётся:
точно установлено, где находится недостающая структура и почему ни trace, ни
projective flux, ни intrinsic cohomology её не заменяют.

Машинные протоколы сохранены в
\texttt{s2t\_family\_affine\_incidence\_exhaustive\_audit.py},
\texttt{s2t\_family\_affine\_incidence\_exhaustive\_results.json},
\texttt{s2t\_even\_projective\_carrier\_exhaustive\_audit.py} и
\texttt{s2t\_even\_projective\_carrier\_exhaustive\_results.json}, а также в
\texttt{s2t\_topological\_family\_selector\_audit.py} и
\texttt{s2t\_topological\_family\_selector\_results.json}, а также в
\texttt{s2t\_spin\_menu\_dirac\_determinant\_line\_audit.py} и
\texttt{s2t\_spin\_menu\_dirac\_determinant\_line\_results.json}.

\end{document}